\documentclass[11pt]{article}
\usepackage{amsmath,amssymb,amsthm}
\usepackage[margin=1in]{geometry}
\usepackage{array}
\usepackage{enumitem}
\usepackage{hyperref}

\newcommand{\N}{\mathbf{N}}
\newcommand{\M}{\mathbf{M}}
\newcommand{\PP}{\mathbf{P}}
\newcommand{\QQ}{\mathbf{Q}}
\newcommand{\JJ}{\mathbf{J}}
\newcommand{\I}{\mathbf{I}}

\title{The $3\times3$ Magic Square of Squares:\\ Structure, Conditions, Eigenmatrix, and Value Limits}
\author{Shaimaa Soltan}
\date{August 2026}

\begin{document}
\maketitle

\begin{abstract}
A summary of the algebraic structure of the $3\times3$ magic square, reduced to
its two shape parameters, together with the conditions for nine distinct cells,
the modular sieves on the offsets, the eigen-structure of the offset matrix
$\N$, and the numerical limits on all the parameters. The governing object is a
single shape ratio $r=C/A$ and an offset matrix $\N=a\PP+c\QQ$.
\end{abstract}

\section{Setup and parametrization}

A $3\times3$ magic square is written row-major as
\[
\begin{pmatrix} A & B & C\\ D & E & F\\ G & H & I\end{pmatrix},
\qquad
\text{all three rows, three columns and both diagonals equal the sum } S.
\]

\paragraph{Center law.} Every line through the center forces
\[
\boxed{\,E=\tfrac{S}{3}\,}\qquad\Longleftrightarrow\qquad S=3E,
\]
because the two cells flanking the center are $E\mp(A-C)$, whose offsets cancel.
Equivalently, each cell plus its point-reflection through the center satisfies
\[
A+I=B+H=C+G=D+F=2E \quad(\text{opposite cells sum to }2E).
\]

\paragraph{Generator basis.} Subtracting the center and factoring gives the
\emph{offset matrix}
\[
\N=a\,\PP+c\,\QQ,\qquad
\PP=\begin{pmatrix}1&-1&0\\-1&0&1\\0&1&-1\end{pmatrix},\quad
\QQ=\begin{pmatrix}0&-1&1\\1&0&-1\\-1&1&0\end{pmatrix},
\]
\[
\N=\begin{pmatrix} a & -a-c & c\\ -a+c & 0 & a-c\\ -c & a+c & -a\end{pmatrix},
\]
where $\PP=\PP^{\mathsf T}$ is symmetric (the ``$a$'' generator) and
$\QQ=-\QQ^{\mathsf T}$ is skew-symmetric (the ``$c$'' generator). Each of
$\PP,\QQ$ has every line summing to $0$.

\paragraph{Full square.} With $\JJ$ the all-ones matrix and center scale $E$,
\[
\boxed{\ \text{Square}=E\,(\JJ+\N)\ },\qquad
\text{cell}=E\,(1+n),\quad n\in\{0,\pm a,\pm c,\pm(a+c),\pm(a-c)\}.
\]

\section{Magic conditions and degrees of freedom}

\begin{center}
\begin{tabular}{lll}
\hline
Object & Independent line-equations & Degrees of freedom\\
\hline
Full magic square & $7$ (of $8$; one is dependent) & $9-7=2$ (plus scale $E$)\\
Semi-magic (rows+cols) & $5$ & $9-5=4$\\
\hline
\end{tabular}
\end{center}

The single dependency among the eight lines is the grand total:
$R_1+R_2+R_3=\text{(all cells)}=C_1+C_2+C_3$. In the normalized variables the two
non-trivial magic conditions are
\[
a+b+c=0,\qquad a+d-c=0
\qquad\Longrightarrow\qquad b=-a-c,\ \ d=c-a,
\]
leaving exactly the two free parameters $a$ and $c$.

\section{Normalized forms}

\[
\underbrace{\text{cells}}_{\text{sum }3E}
\ \xrightarrow{\ \div E\ }\
\M=\JJ+\N \quad(\text{center }1,\ \text{sum }3)
\ \xrightarrow{\ \div 3\ }\
\tfrac13\M \quad(\text{center }\tfrac13,\ \text{sum }1).
\]
\[
\M=\begin{pmatrix}
1+a & 1-a-c & 1+c\\
1-a+c & 1 & 1+a-c\\
1-c & 1+a+c & 1-a
\end{pmatrix}.
\]
\begin{itemize}[nosep]
\item \textbf{Center-1 form} ($\div E$): \emph{square-friendly}. Since $E=\rho^2$
is a perfect square, an integer-square cell $\Leftrightarrow$ a rational-square
entry of $\M$. Use this frame for the square hunt.
\item \textbf{Sum-1 form} ($\div 3E$): \emph{barycentric}. Center $=\tfrac13$,
each flanking pair $=\tfrac23$. Do \emph{not} use for squares: dividing by $3$
(not a square) breaks the integer$\leftrightarrow$rational-square correspondence.
\end{itemize}

\section{Reduced problem: shape ratio and forbidden values}

Factor $a$ from $\N$ and set the shape ratio
\[
\boxed{\,r=\frac{c}{a}=\frac{C}{A}\,},\qquad \N=a\,(\PP+r\,\QQ).
\]
The four offset magnitudes are proportional to $1,\ r,\ 1+r,\ 1-r$
(directions $0^\circ,90^\circ,45^\circ,135^\circ$ in the $(a,c)$ plane).

\paragraph{Symmetry.} The $8$ magic-square symmetries act on the ratio as
$r\mapsto\{-r,\ 1/r,\ -1/r\}$, so every distinct shape has a representative in
\[
\boxed{\,0<r<1\,}\quad(\text{fundamental domain}).
\]

\paragraph{Forbidden ratios (no $9$ distinct cells).} A cell duplicates or hits
the center exactly when
\[
\boxed{\,r\in\{0,\ \pm\tfrac12,\ \pm1,\ \pm2,\ \infty\}\,}
\]
i.e.\ the axes ($0,\infty$), the diagonals ($\pm1$), and their $45^\circ$
reflection partners ($\pm\tfrac12,\pm2$). In the fundamental range the only
forbidden value is $r=\tfrac12$ (endpoints $0,1$ excluded).

\paragraph{Positivity (all nine cells $>0$).}
\[
|a|<\frac{1}{\max\bigl(1,|r|,|1+r|,|1-r|\bigr)}=\frac{1}{1+r}\quad(0<r<1).
\]
This bounds the offset scale $a$; it does \emph{not} bound $r$.

\paragraph{Two kinds of boundary (radius vs.\ direction).} Read $(r,a)$ as polar
coordinates: $r$ is the \emph{direction} and $a$ is the \emph{radius} along it.
The two limits act on different coordinates:
\begin{itemize}[nosep]
\item The \textbf{forbidden ratios} delete whole \emph{directions} (rays through
the origin): $r\in\{0,\pm\tfrac12,\pm1,\pm2,\infty\}$. This is a limit \emph{on
$r$}---it removes values of the shape.
\item \textbf{Positivity} caps the \emph{radius}: $|a|<1/(1+r)$. This is a limit
\emph{on $a$}. Its numerical value is a \emph{function of} $r$, but it never
deletes a direction: every admissible $r\in(0,1)$ keeps a nonempty window
$|a|<1/(1+r)$. A direction-dependent cap on the radius is not a restriction on the
direction.
\end{itemize}
Thus ``bounds $a$, not $r$'' means exactly: positivity excludes no value of $r$;
it only fixes how large $|a|$ may be once the direction $r$ is chosen. The tightest
cell is $1-a(1+r)$ (largest magnitude $1+r$ on $0<r<1$), which is why the cap is
$1/(1+r)$.

\medskip
\noindent\textbf{Geometry in the $(a,c)$-plane.} Positivity is the open hexagon
about the origin
\[
\max\bigl(|a|,\ |c|,\ |a+c|,\ |a-c|\bigr)<1,
\]
and the forbidden ratios are the finite set of \emph{rays} $c=r\,a$,
$r\in\{0,\pm\tfrac12,\pm1,\pm2,\infty\}$, deleted from it. The hexagon caps the
radius (the $a$-bound, tighter in some directions); the deleted rays cut out
directions (the $r$-limit)---two different boundaries on the same $2$-dimensional
shape region.

\section{Square conditions and the elliptic curve}

Nine cells are perfect squares iff, in the center-$1$ frame,
\[
1\pm a,\quad 1\pm c,\quad 1\pm(a+c),\quad 1\pm(a-c)
\]
are all rational squares (the center $1$ is free). Each pair
``$1\pm x$ both squares'' is a rational point on the conic $p^2+q^2=2$ (genus $0$,
$\cong$ a line). The two \emph{coupled} conditions on $a+c$ and $a-c$ intersect
these into a \textbf{genus-$1$ elliptic curve}; its rational points are the
many-square magic squares.

\medskip
\noindent\textbf{Coordinates: scale versus shape.} The two free parameters are
$(a,c)$; equivalently use $(a,r)$ with
\[
r=\frac{C}{A}=\frac{c}{a}\qquad(\text{scale-invariant shape}),\qquad
a\ (\text{scale}),\qquad c=a\,r .
\]
Here $r$ is \emph{built from} $a$ and $c$, but $(a,r)$ is a genuine coordinate
system: $r$ is unchanged under a common rescaling $(a,c)\mapsto(\lambda a,\lambda c)$,
while $a$ carries the magnitude.

\medskip
\noindent\textbf{Independence of the two conditions (not of the variables).}
\begin{itemize}[nosep]
\item The \emph{distinctness} condition (avoid the forbidden ratios) depends
\emph{only on the shape} $r$; scaling $a$ leaves it untouched.
\item The \emph{squareness} conditions depend on \emph{both} $a$ and $r$
(the actual values $1\pm a,\,1\pm ar,\dots$ must be rational squares).
\end{itemize}
Hence the two \emph{conditions} are logically independent---choosing a permitted
shape neither implies nor precludes the existence of a scale $a$ making the cells
square---so they cannot contradict each other. In particular, squaring a
magnitude never forces $r$ into the forbidden set: the record family below squares
one diagonal magnitude fully at $r\notin$ forbidden set.

\section{Modular sieves on the offsets}

A square is $0$ or $1\pmod 3$, $0,1,4\pmod 8$, etc. Writing the nine cells as
$E\pm(\text{offset})$ gives necessary conditions on $A,C$:
\[
\begin{array}{lll}
\text{mod }3: & A\equiv C\equiv 0 & (\text{necessary for } \ge 7 \text{ squares})\\
\text{mod }8/16: & A\equiv C\equiv 0 \bmod 8 &\\
\text{mod }5\ (\text{last digit}): & A\equiv 0 \text{ or } C\equiv 0 &\\
\hline
\text{mod }240=16\cdot3\cdot5: & \text{only } 1/2880 \text{ of }(A,C)\text{ survive} &
\end{array}
\]
Every record family passes all sieves. A pair value from
$\mathrm{squarepairs}(2E)$ is already a perfect square, hence passes every
sieve automatically---so drawing free cells from square-pairs is strictly
stronger than the sieve and runs in $O(\text{pairs}^2)$.

\section{Eigenmatrix of $\N$}

The offset matrix is rank $2$ and generates its own powers:
\[
\boxed{\ \N^2=(a^2-c^2)\,(3\I-\JJ)\ },\qquad
\boxed{\ \N^3=3(a^2-c^2)\,\N\ }.
\]
Writing $\kappa=a^2-c^2=(a+c)(a-c)$ (the product of the two diagonal magnitudes),
the eigenvalues are
\[
\operatorname{spec}(\N)=\Bigl\{\,0,\ +\sqrt{3\kappa},\ -\sqrt{3\kappa}\,\Bigr\}.
\]
The all-ones vector $\mathbf 1$ is the $0$-eigenvector (line sums vanish). Note
$\N^2$ has only two distinct entries---squaring collapses the shape to the single
scalar $\kappa$; the nine-distinct structure lives only in the \emph{linear}
$\N$.

\subsection*{Derivation of $\N^{1/2}$}

Since $\N$ generates its own powers
\[
\N^2=\kappa\,(3\I-\JJ),\qquad \N^3=3\kappa\,\N,\qquad
\N^4=\N\cdot\N^3=3\kappa\,\N^2=3\kappa^2(3\I-\JJ),\qquad \kappa=a^2-c^2,
\]
any polynomial in $\N$ closes on the two building blocks $\N$ and $(3\I-\JJ)$.
Because $\N$ annihilates the all-ones vector ($0$ eigenvalue), the square root has
no constant term, so it is a \emph{quadratic in $\N$}:
\[
\N^{1/2}=q\,\N+p\,\N^2 .
\]
Squaring and reducing with the identities above,
\[
\bigl(q\N+p\N^2\bigr)^2
= q^2\N^2+2pq\,\N^3+p^2\N^4
= \underbrace{\bigl(q^2\kappa+3p^2\kappa^2\bigr)}_{\text{coeff. of }(3\I-\JJ)}(3\I-\JJ)
\;+\;\underbrace{6pq\kappa}_{\text{coeff. of }\N}\,\N .
\]
Matching this to $\N$ (no $(3\I-\JJ)$ part, unit $\N$ part) gives the two
coefficient equations
\[
6pq\kappa=1,\qquad q^2\kappa+3p^2\kappa^2=0
\quad\Longrightarrow\quad
pq=\frac{1}{6\kappa},\qquad q^2=-3\kappa\,p^2 .
\]
Eliminating in turn,
\[
\boxed{\ \N^{1/2}=q\,\N+p\,\N^2\ },\qquad
\boxed{\ q^4=\frac{-1}{12\kappa},\qquad p^4=\frac{-1}{108\,\kappa^3}\ },
\qquad \kappa=a^2-c^2 .
\]
Because $p^4,q^4<0$, $\N^{1/2}$ is \textbf{complex} (forced by the eigenvalue
$-\sqrt{3\kappa}$). It is \emph{semi-magic} (rows and columns sum to $0$) but not
full magic (diagonals differ), and it fails to exist when $\kappa=0$, i.e.\ at the
forbidden diagonal ratios $r=\pm1$ (there $\N$ is nilpotent, $\N^2=0$).

\subsection*{Limits on the coefficients $p,q$ from the $a,c$ limits}

Everything passes through the single invariant $\kappa=a^2-c^2=a^2(1-r^2)$. The
shape and positivity limits of \S4 pin its range:
\[
0<r<1\ \Rightarrow\ \kappa>0,\qquad
|a|<\frac{1}{1+r}\ \Rightarrow\ \kappa<\frac{1-r}{1+r}<1
\qquad\Longrightarrow\qquad
\boxed{\,0<\kappa<1\,}.
\]
The supremum $\kappa\to1$ is approached as $r\to0,\ a\to1$; the infimum
$\kappa\to0$ as $r\to1$ (the degenerate diagonal). From
$|p|^4=\tfrac{1}{108\kappa^3}$ and $|q|^4=\tfrac{1}{12\kappa}$,
\[
|p|=\frac{1}{\bigl(108\,\kappa^3\bigr)^{1/4}},\qquad
|q|=\frac{1}{\bigl(12\,\kappa\bigr)^{1/4}},
\]
which, on $0<\kappa<1$, give
\[
\boxed{\ |p|>108^{-1/4}\approx0.3102,\qquad |q|>12^{-1/4}\approx0.5373\ }
\qquad(\text{both}\ \to\infty\ \text{as}\ \kappa\to0,\ r\to1).
\]
The two coefficients are tied by clean invariants:
\[
|p\,q|=\frac{1}{6\kappa}>\frac16,\qquad
\left|\frac{q}{p}\right|=\sqrt{3\kappa}<\sqrt3
\quad\bigl(=\ |\mu|,\ \text{the nonzero eigenvalue}\bigr).
\]
So the square-root coefficients have \emph{lower} bounds $\bigl(108^{-1/4},\,
12^{-1/4}\bigr)$ at the maximal shape ($\kappa\to1$) and \emph{blow up} at the
degenerate diagonal ($\kappa\to0$, $r\to\pm1$), consistent with $\N^{1/2}$ ceasing
to exist there. The phase of $p,q$ is fixed by $p^4,q^4<0$ (fourth roots of a
negative real, i.e.\ arguments $\tfrac{\pi}{4}+\tfrac{k\pi}{2}$).

\subsection*{Remark: factoring $6pq$ yields the diagonal magnitudes}

Because $\kappa=(a+c)(a-c)$ and $pq=\tfrac{1}{6\kappa}$, the reciprocal
\[
6pq=\frac{1}{\kappa}=\frac{1}{(a+c)(a-c)}
\]
carries the whole diagonal structure. A single value $pq$ fixes only the invariant
$\kappa$ (a hyperbola $a^2-c^2=\text{const}$, with the shape $r$ still free), but a
\emph{factorization} $6pq=m\cdot n$ distributes it into the two diagonal
magnitudes:
\[
a+c=\frac1m,\qquad a-c=\frac1n
\quad\Longrightarrow\quad
a=\frac{m+n}{2mn},\quad c=\frac{n-m}{2mn},\quad r=\frac{c}{a}=\frac{n-m}{n+m}.
\]
Each factor pair is a distinct rational shape. For example $pq=35$ gives
$6pq=210=2\cdot3\cdot5\cdot7$; the pairing with the $6$ is $(m,n)=(6,35)$, whence
\[
a+c=\tfrac16,\quad a-c=\tfrac1{35},\quad a=\tfrac{41}{420},\quad c=\tfrac{29}{420},
\quad r=\tfrac{29}{41},
\]
and the eight factor pairs of $210$ give eight admissible shapes (all in $(0,1)$,
none equal to $\tfrac12$). Thus the \emph{product} $pq$ sets the invariant, while
\emph{factoring} $6pq$ is exactly factoring $(a+c)(a-c)$ into the two diagonals.

\section{Shape $\leftrightarrow$ factorization}

The additive representation $a+c=\tfrac1m,\ a-c=\tfrac1n$ (with $mn=6pq$) links the
\emph{shape} to the \emph{factorization} through two exact, additive bridges.

\subsection*{The shape numerators are symmetric about $n$}

Over the common denominator $2mn$, the offsets are
\[
a=\frac{m+n}{2mn},\qquad c=\frac{n-m}{2mn},
\]
so the two numerators are $n+m$ and $n-m$: symmetric about $n$ with half-gap $m$,
\[
c\text{-num}=n-m,\qquad a\text{-num}=n+m,\qquad
\text{center}=\frac{(n+m)+(n-m)}{2}=n,\qquad \text{half-gap}=m.
\]
The Möbius map on the shape recovers the factor pair exactly:
\[
\boxed{\ \frac{1+r}{1-r}=\frac{n}{m}\ },\qquad r=\frac{c}{a}=\frac{n-m}{n+m},
\]
and the reciprocal shape $r\mapsto1/r$ flips the sign, $\tfrac{1+1/r}{1-1/r}=-\tfrac nm$.
(This is the same $45^\circ$ Möbius map of \S5, turning axis magnitudes into
diagonal ones.) In the special case $m=6$ the center is $n=pq$ itself and the
half-gap is $6$, e.g.\ $71=77-6$, $83=77+6$, giving $\tfrac{1+r}{1-r}=\tfrac{77}{6}$.

\subsection*{Euler: two sums of two squares factor the center}

The magic-square richness lives in $2c=u^2+v^2$; a center with two \emph{distinct}
such representations factors by Euler's method:
\[
N=a^2+b^2=c^2+d^2\quad(\{a,b\}\neq\{c,d\})
\;\Longrightarrow\;
\gcd(ad\pm bc,\ N)\ \text{are nontrivial factors of }N.
\]
For instance $65=1^2+8^2=4^2+7^2\Rightarrow\gcd(1\cdot7\mp8\cdot4,65)=5,13$, and
$2501=1^2+50^2=10^2+49^2\Rightarrow41,61$. Hence
\[
\text{many additive reps of }2c
\iff c\text{ built from primes }\equiv1\!\!\pmod4
\iff \text{Euler factors }c .
\]
Richness (\,$\mathrm{square\_pairs}(2c)$\,) and factorability are the same fact:
this is why the record families sit at highly composite centers.

\subsection*{Why neither factors ``for free''}

Both bridges are genuine identities, yet neither is a shortcut, for the same
reason:
\begin{itemize}[nosep]
\item $\tfrac{1+r}{1-r}=\tfrac nm$ is a \emph{bijection} shape $\leftrightarrow$
factor pair. Using it presupposes the shape $r=C/A$, which is itself built from
$m,n$; it returns what was put in.
\item Euler's $\gcd$ is instant \emph{given} two representations, but finding even
one (searching $a$ with $N-a^2$ square) costs $\Theta(\sqrt N)$.
\end{itemize}
Any fixed arithmetic rule that output a nontrivial factor of $N$ would either take
a factor as (possibly hidden) input or scan $\Theta(\sqrt N)$ candidates with a
divisibility test---so no residue/floor/Möbius expression escapes the
$\sqrt N$ cost. The additive lens \emph{organizes} the problem correctly; it does
not cheapen the factoring.

\subsection*{Where ``peel 2, 3'' legitimately leads: the $6k\pm1$ wheel}

There \emph{is} a real, correct algorithm at the end of this line, and it is the
one the $6k\pm1$ structure points to. Consider the candidate modulus
\[
M=(N\bmod 6)+6m,\qquad m=0,1,2,\dots
\]
For $N$ coprime to $6$ (i.e.\ after peeling the factors $2$ and $3$), $N\bmod6\in\{1,5\}$,
so $M$ runs through one residue class $\{5,11,17,23,\dots\}$ or $\{1,7,13,19,\dots\}$
---exactly the $6k\pm1$ numbers, the only possible prime factors of $N$. Testing
$N\bmod M=0$ is a divisibility check, and scanning $m$ is
\textbf{wheel ($6k\pm1$) trial division}: a genuine method, $\sim\!3\times$ faster
than testing every integer because multiples of $2$ and $3$ are skipped. Two
requirements make it correct, and one caveat bounds it:
\begin{itemize}[nosep]
\item \textbf{Scan $m$.} A single fixed $m$ tests one candidate; factoring needs
$m$ incremented until $M\mid N$ or $M>\sqrt N$.
\item \textbf{Both classes.} A factor may lie in \emph{either} $\{1,5\}\bmod6$.
E.g.\ $55=5\cdot11$ has $N\equiv1\pmod6$ yet both primes are $\equiv5$; scanning
only the class $N\bmod6$ misses them.
\item \textbf{Still $\Theta(\sqrt N)$.} The wheel removes a constant fraction of
candidates but marches $M$ up to $\sqrt N$; it is a constant-factor win, not an
asymptotic one.
\end{itemize}
This is the honest terminus of the ``peel $2,3$, then everything is $6k\pm1$''
structure (the same $6k\pm1$ blocks $5=6-1$, $7=6+1$ that seed the flanking-pair
arithmetic): the residue and M\"obius shortcuts dissolve under a fair sample of
semiprimes, but the $6k\pm1$ wheel survives---because it \emph{is} trial division,
in its $2,3$-peeled form, and therefore respects the $\sqrt N$ cost rather than
evading it.

\subsection*{Closing the loop: the $\{1,2,3,+,-,*\}$ canonical form}

The same $6k\pm1$ fact gives a canonical representation of \emph{every} natural
number over the alphabet $\{1,2,3,+,-,*\}$. Factor $N$ with the wheel; then $2$ and
$3$ are atoms, and every prime $p>3$ is
\[
p = 2\cdot3\cdot k \pm 1,\qquad k=\frac{p\mp1}{6},
\]
with $k$ written in the same form, recursively. Thus
\[
N=\prod_i p_i,\qquad p_i=2\cdot3\cdot(\text{canonical }k_i)\pm1 .
\]
For example
\[
5=(2{\cdot}3{\cdot}1-1),\quad 7=(2{\cdot}3{\cdot}1+1),\quad 17=(2{\cdot}3{\cdot}3-1),
\]
\[
425=5\cdot5\cdot17=(2{\cdot}3{\cdot}1-1)(2{\cdot}3{\cdot}1-1)(2{\cdot}3{\cdot}3-1),
\qquad
180625=425^2=(2{\cdot}3{\cdot}1-1)^4(2{\cdot}3{\cdot}3-1)^2,
\]
and primes recurse, e.g.\ $2501=41\cdot61=(2{\cdot}3{\cdot}(2{\cdot}3{\cdot}1+1)-1)\,
(2{\cdot}3{\cdot}2{\cdot}(2{\cdot}3{\cdot}1-1)+1)$.
The wheel is the engine (it yields the primes); the $6k\pm1$ identity is the rule
that rebuilds each prime from $\{1,2,3,\pm1,*\}$. This is the same single fact---%
\emph{after $2$ and $3$, everything is $6k\pm1$}---that underlies the flanking-pair
richness ($2c=u^2+v^2$ needs primes $\equiv1\bmod4$ but the blocks are $6k\pm1$),
the factorization wheel, and this canonical form: one structure, three faces.

\section{Summary of value limits}

\begin{center}
\renewcommand{\arraystretch}{1.35}
\begin{tabular}{p{4.7cm} p{9.5cm}}
\hline
\textbf{Quantity} & \textbf{Condition / limit}\\
\hline
Center $E$ & $E=\rho^2$ (perfect square) so cell$/E$ is a rational square; sum $=3E$.\\
Degrees of freedom & $2$ shape (a,c) $+\,1$ scale $E$; semi-magic has $4$.\\
Shape ratio $r=C/A$ & fundamental domain $0<r<1$; symmetry $r\sim\{-r,1/r,-1/r\}$.\\
Forbidden $r$ (no $9$ distinct) & $r\in\{0,\pm\tfrac12,\pm1,\pm2,\infty\}$ (only $\tfrac12$ inside $(0,1)$).\\
Offset scale $a$ (positivity) & $|a|<\dfrac{1}{1+r}$ for $0<r<1$.\\
Offsets $A,C$ (mod sieve, $\ge7$ sq.) & $A\equiv C\equiv0\ (\mathrm{mod}\ 3)$; combined mod $240$ keeps $1/2880$.\\
Eigenvalues of $\N$ & $0,\ \pm\sqrt{3(a^2-c^2)}$; real iff $|r|<1$.\\
Invariant $\kappa$ & $\kappa=a^2-c^2=(a+c)(a-c)$; $\kappa>0$ on $0<r<1$; $\kappa=0$ at $r=\pm1$.\\
Cells vs.\ sum & every cell $<2E=\tfrac23 S$ (opposite $>0$).\\
\hline
\end{tabular}
\end{center}

\section{Worked examples: the $425$ and $195364$ squares}

In each grid the nine cells are shown with perfect squares marked $^\star$ (root
given in the caption). The right-hand column lists the three \textbf{row sums};
the bottom row lists the three \textbf{column sums}; and the bottom-right corner
gives the two \textbf{diagonal sums} ($\diagdown$ = main, $\diagup$ = anti).

\begin{table}[h]
\centering
\begin{tabular}{|rrr||r|}
\hline
$139129^\star$ & $83521^\star$ & $319225^\star$ & $\mathbf{541875}$\\
$360721$ & $180625^\star$ & $529^\star$ & $\mathbf{541875}$\\
$42025^\star$ & $277729^\star$ & $222121$ & $\mathbf{541875}$\\
\hline\hline
$\mathbf{541875}$ & $\mathbf{541875}$ & $\mathbf{541875}$
   & \shortstack{$\diagdown\,\mathbf{541875}$\\[2pt]$\diagup\,\mathbf{541875}$}\\
\hline
\end{tabular}
\caption{\textbf{$425$ family} --- root $\rho=425$, center $E=425^2=180625$, magic
sum $3E=541875$. All rows, columns \emph{and} both diagonals equal $541875$.
\textbf{Magic}, $7$ squares: $373^2,289^2,565^2,425^2,23^2,205^2,527^2$; the two
non-squares are $360721$ and $222121$.}
\end{table}

\begin{table}[h]
\centering
\begin{tabular}{|rrr||r|}
\hline
$40804^\star$ & $386884^\star$ & $158404^\star$ & $\mathbf{586092}$\\
$312964$ & $195364^\star$ & $77764$ & $\mathbf{586092}$\\
$232324^\star$ & $3844^\star$ & $349924$ & $\mathbf{586092}$\\
\hline\hline
$\mathbf{586092}$ & $\mathbf{586092}$ & $\mathbf{586092}$
   & \shortstack{$\diagdown\,\mathbf{586092}$\\[2pt]$\diagup\,\mathbf{586092}$}\\
\hline
\end{tabular}
\caption{\textbf{$195364$ magic square} --- center $E=442^2=195364$, magic sum
$3E=586092$. All rows, columns and both diagonals equal $586092$. \textbf{Magic},
$6$ squares:
$202^2,622^2,398^2,442^2,482^2,62^2$; non-squares $312964,77764,349924$. This is the
best \emph{magic} square-count reachable at this center.}
\end{table}

\begin{table}[h]
\centering
\begin{tabular}{|rrr||r|}
\hline
$3844^\star$ & $33124^\star$ & $56644^\star$ & $93612$\\
$158404^\star$ & $195364^\star$ & $232324^\star$ & $\mathbf{586092}$\\
$334084^\star$ & $357604^\star$ & $386884^\star$ & $1078572$\\
\hline\hline
$496332$ & $\mathbf{586092}$ & $675852$
   & \shortstack{$\diagdown\,\mathbf{586092}$\\[2pt]$\diagup\,\mathbf{586092}$}\\
\hline
\end{tabular}
\caption{\textbf{$195364$ nine-square} (point-symmetric) --- all $9$ cells are
squares ($62^2,182^2,238^2,398^2,442^2,482^2,578^2,598^2,622^2$), but \textbf{NOT
magic}. Only the center cross reaches $586092$: middle row, middle column, and
\emph{both diagonals} ($\diagdown=\diagup=586092$, shown in the corner). The outer
rows ($93612,\,1078572$) and outer columns ($496332,\,675852$) are off by $\pm X$
($X=492480$ for rows) --- the signature of a point-symmetric grid that is not
magic.}
\end{table}

\begin{table}[h]
\centering
\begin{tabular}{|rrr||r|}
\hline
$349924$ & $33124^\star$ & $386884^\star$ & $769932$\\
$232324^\star$ & $195364^\star$ & $158404^\star$ & $\mathbf{586092}$\\
$3844^\star$ & $357604^\star$ & $40804^\star$ & $402252$\\
\hline\hline
$\mathbf{586092}$ & $\mathbf{586092}$ & $\mathbf{586092}$
   & \shortstack{$\diagdown\,\mathbf{586092}$\\[2pt]$\diagup\,\mathbf{586092}$}\\
\hline
\end{tabular}
\caption{\textbf{$195364$ eight-square} (point-symmetric) --- $8$ of the $9$ cells
are squares ($182^2,622^2,482^2,442^2,398^2,62^2,598^2,202^2$); only the top-left
$349924$ is not ($\sqrt{349924}\approx591.54$). Here \emph{all three columns} and
\emph{both diagonals} \emph{and} the middle row equal $586092$ --- six of the eight
lines are magic --- but the top and bottom rows are off by $\pm183840$
($769932$ and $402252$). It is the closest approach to a magic $8$-square: to make
it magic the top-left cell is forced to $586092-33124-386884=166084=407.5^2\ldots$,
which is \emph{not} a perfect square. That single stubborn cell is the $7\!\to\!8$
wall in one picture.}
\end{table}

\section{Records and open problem}

\begin{itemize}[nosep]
\item \textbf{$9$ distinct squares, non-magic:} easy---any center with
$\ge 4$ flanking square-pairs (point-symmetric grid).
\item \textbf{Magic, $7$ squares:} the ``$425$ family'' at root $\rho=425$,
center $E=180625$, shape $r=-247/825$ (in $(0,1)$ up to symmetry, avoids
$\tfrac12$). Two of the four magnitudes are fully squared.
\item \textbf{Magic, $8$ or $9$ squares:} \emph{open}. Requires all four
magnitudes $1\pm\{a,c,a+c,a-c\}$ simultaneously square---a rational point on the
elliptic curve. No local/modular obstruction rules out any admissible $r$; the
obstruction, if any, is global.
\end{itemize}

\section*{Acknowledgments}

This note grew out of an extended exploration by the author, reducing the
$3\times3$ magic square of squares to its two-parameter shape $(a,c)$, the single
scale-invariant ratio $r=C/A$, and the offset matrix $\N=a\PP+c\QQ$. The symbolic
identities ($\N^2=(a^2-c^2)(3\I-\JJ)$, $\N^3=3(a^2-c^2)\N$, the closed form of
$\N^{1/2}$), the modular sieves, and the numerical checks throughout were verified
computationally with \texttt{SymPy} and \texttt{NumPy}. The author thanks these
open-source tools, whose symbolic and numerical support made the verification of
every claim in this note routine.

\end{document}
